\documentclass[11pt,a4paper]{article}

\usepackage{fontspec}
\usepackage{polyglossia}
\usepackage{geometry}
\usepackage{graphicx}
\usepackage{xcolor}
\usepackage{booktabs}
\usepackage{array}
\usepackage{tabularx}
\usepackage{titlesec}
\usepackage{enumitem}
\usepackage{hyperref}
\usepackage{caption}
\usepackage{fancyhdr}

\setmainlanguage{english}
\setmainfont{Arial}

\geometry{top=20mm,bottom=20mm,left=18mm,right=18mm}
\setlength{\headheight}{14pt}
\graphicspath{{customer_product_analysis/charts/}}

\definecolor{ink}{HTML}{17202A}
\definecolor{muted}{HTML}{5D6D7E}
\definecolor{blue}{HTML}{2D67B4}
\definecolor{green}{HTML}{238C5F}
\definecolor{red}{HTML}{C44348}
\definecolor{panel}{HTML}{F4F6F8}

\hypersetup{
    colorlinks=true,
    linkcolor=blue,
    urlcolor=blue
}

\pagestyle{fancy}
\fancyhf{}
\lhead{\textcolor{muted}{Customer and Product Profitability Analysis}}
\rhead{\textcolor{muted}{Global SuperStore}}
\cfoot{\textcolor{muted}{\thepage}}
\renewcommand{\headrulewidth}{0pt}

\titleformat{\section}
  {\Large\bfseries\color{ink}}
  {}
  {0pt}
  {}

\titleformat{\subsection}
  {\large\bfseries\color{ink}}
  {}
  {0pt}
  {}

\setlist[itemize]{leftmargin=*, itemsep=2pt, topsep=4pt}
\setlength{\parindent}{0pt}
\setlength{\parskip}{6pt}
\captionsetup{font=small,labelfont=bf}

\newcommand{\metric}[2]{
    \begin{minipage}[t]{0.18\textwidth}
        \colorbox{panel}{
            \begin{minipage}[t][25mm][c]{0.92\linewidth}
                \centering
                {\Large\bfseries\color{ink} #1}\\[2mm]
                {\small\color{muted} #2}
            \end{minipage}
        }
    \end{minipage}
}

\begin{document}

\begin{titlepage}
    \thispagestyle{fancy}
    \vspace*{8mm}
    {\Huge\bfseries\color{ink} Customer and Product Profitability Analysis\par}
    \vspace{3mm}
    {\Large\color{muted} A deep dive into profit concentration, customer value, and product-level losses\par}
    \vspace{12mm}

    \metric{795}{Customers}
    \hfill
    \metric{3,788}{Products}
    \hfill
    \metric{427}{Customers for 80\% profit}
    \hfill
    \metric{67}{Loss-making customers}
    \hfill
    \metric{3}{Segments}

    \vspace{14mm}
    \section*{Executive Summary}
    This analysis extends the SuperStore sales performance work by focusing on customer and product profitability.
    The key finding is that profit is meaningful across a broad customer base rather than being dominated by only a few accounts.
    At the same time, losses are clearly concentrated in a smaller group of customers and specific products, especially 3D printers, machines, and tables.

    \vspace{4mm}
    \begin{tabularx}{\textwidth}{>{\bfseries}p{45mm}X}
        \toprule
        Business question & Which customers and products create the most profit, which ones destroy margin, and how concentrated is profit across the customer base? \\
        Scope & 795 customers, 3,788 products, 2011--2014 \\
        Main customer finding & 427 profitable customers generate 80\% of positive customer profit \\
        Main product finding & Losses are concentrated in selected Technology Machines and Furniture Tables products \\
        Main recommendation & Manage profitability at customer and product level, not by sales volume alone \\
        \bottomrule
    \end{tabularx}

    \vfill
    {\color{muted} Prepared as a data analytics portfolio case study.}
\end{titlepage}
\setcounter{page}{2}

\section{1. Customer Profit Concentration}

The customer base contains 795 unique customers. The top 427 profitable customers generate 80\% of positive customer profit.
This is a relatively broad profit base, which means the business is not dependent on only a handful of accounts.

\begin{figure}[h]
    \centering
    \includegraphics[width=\textwidth]{customer_profit_pareto.png}
    \caption{Customer profit Pareto curve}
\end{figure}

\begin{tabularx}{\textwidth}{lrrr}
    \toprule
    Class & Customers & Sales & Profit \\
    \midrule
    A & 427 & \$7,589,112 & \$1,227,368 \\
    B & 169 & \$2,343,587 & \$229,796 \\
    C & 132 & \$1,761,355 & \$77,384 \\
    Loss & 67 & \$948,851 & -\$65,513 \\
    \bottomrule
\end{tabularx}

\section{2. Most Profitable Customers}

The most profitable customers combine meaningful sales volume with strong margin performance.
They are not simply the largest buyers; they also generate attractive profit per dollar of sales.

\begin{figure}[h]
    \centering
    \includegraphics[width=\textwidth]{top_customers_profit.png}
    \caption{Top customers by profit}
\end{figure}

\begin{tabularx}{\textwidth}{lrrr}
    \toprule
    Customer & Profit & Sales & Profit margin \\
    \midrule
    Tamara Chand & \$8,673 & \$37,453 & 23.2\% \\
    Raymond Buch & \$8,453 & \$29,601 & 28.6\% \\
    Sanjit Chand & \$8,205 & \$26,524 & 30.9\% \\
    Hunter Lopez & \$7,817 & \$30,246 & 25.8\% \\
    Bill Eplett & \$7,410 & \$28,480 & 26.0\% \\
    \bottomrule
\end{tabularx}

\section{3. Loss-Making Customers}

There are 67 customers with negative total profit.
These customers are not automatically bad accounts, but they should be reviewed before additional discounts or promotions are expanded to them.
The review should focus on discounting, product mix, shipping economics, and whether losses are isolated or recurring.

\begin{figure}[h]
    \centering
    \includegraphics[width=\textwidth]{bottom_customers_profit.png}
    \caption{Bottom customers by profit}
\end{figure}

\begin{tabularx}{\textwidth}{lrrr}
    \toprule
    Customer & Profit & Sales & Profit margin \\
    \midrule
    Cindy Stewart & -\$6,152 & \$13,432 & -45.8\% \\
    Luke Foster & -\$3,644 & \$13,422 & -27.2\% \\
    Grant Thornton & -\$3,578 & \$20,226 & -17.7\% \\
    Candace McMahon & -\$2,799 & \$14,798 & -18.9\% \\
    Skye Norling & -\$2,638 & \$14,225 & -18.5\% \\
    \bottomrule
\end{tabularx}

\section{4. Product Profitability}

The strongest products are mostly Technology items, especially copiers, phones, and routers.
These products combine scale with strong profit margins.

\begin{figure}[h]
    \centering
    \includegraphics[width=\textwidth]{top_products_profit.png}
    \caption{Top products by profit}
\end{figure}

\begin{tabularx}{\textwidth}{Xrrr}
    \toprule
    Product & Profit & Sales & Profit margin \\
    \midrule
    Canon imageCLASS 2200 Advanced Copier & \$25,200 & \$61,600 & 40.9\% \\
    Cisco Smart Phone, Full Size & \$17,239 & \$76,441 & 22.6\% \\
    Motorola Smart Phone, Full Size & \$17,027 & \$73,159 & 23.3\% \\
    Hoover Stove, Red & \$11,808 & \$31,664 & 37.3\% \\
    Sauder Classic Bookcase, Traditional & \$10,672 & \$39,110 & 27.3\% \\
    \bottomrule
\end{tabularx}

\section{5. Product Losses}

The largest product losses are concentrated in selected Technology Machines and Furniture Tables products.
This aligns with the earlier category-level finding that Tables reduce Furniture profitability.
High average discounts are a common pattern among the weakest products.

\begin{figure}[h]
    \centering
    \includegraphics[width=\textwidth]{bottom_products_profit.png}
    \caption{Bottom products by profit}
\end{figure}

\begin{tabularx}{\textwidth}{Xrrr}
    \toprule
    Product & Profit & Sales & Profit margin \\
    \midrule
    Cubify CubeX 3D Printer Double Head Print & -\$8,880 & \$11,100 & -80.0\% \\
    Lexmark MX611dhe Monochrome Laser Printer & -\$4,590 & \$16,830 & -27.3\% \\
    Motorola Smart Phone, Cordless & -\$4,447 & \$38,935 & -11.4\% \\
    Cubify CubeX 3D Printer Triple Head Print & -\$3,840 & \$8,000 & -48.0\% \\
    Bevis Round Table, Adjustable Height & -\$3,650 & \$5,655 & -64.5\% \\
    \bottomrule
\end{tabularx}

\section{6. Segment View}

Customer segment is less important than product mix and discount behavior.
Consumer, Corporate, and Home Office all show similar profit margins around 11.5--12.0\%.

\begin{tabularx}{\textwidth}{lrrr}
    \toprule
    Segment & Sales & Profit & Profit margin \\
    \midrule
    Consumer & \$6,508,141 & \$749,240 & 11.5\% \\
    Corporate & \$3,824,808 & \$442,786 & 11.6\% \\
    Home Office & \$2,309,956 & \$277,009 & 12.0\% \\
    \bottomrule
\end{tabularx}

\section{Recommendations}

\begin{itemize}
    \item Protect high-profit customers with targeted retention and service quality.
    \item Review the 67 loss-making customers before expanding discounts or promotions to them.
    \item Investigate high-loss products, especially 3D printers, machines, and tables.
    \item Use customer profitability, not sales alone, as the primary account-prioritization metric.
    \item Keep segment-level strategy broad, but manage discounting and product mix at customer/product level.
\end{itemize}

\section{Appendix: Reproducibility}

The analysis can be reproduced with:

\begin{verbatim}
python scripts/customer_product_analysis.py
\end{verbatim}

\end{document}
